\documentclass[11pt,a4paper]{ctexart}

% ---------- Page geometry ----------
\usepackage[left=2cm,right=2cm,top=2cm,bottom=2cm]{geometry}

% ---------- Math ----------
\usepackage{amsmath}
\usepackage{amssymb}
\usepackage{amsfonts}
\usepackage{mathtools}
\usepackage{amsthm}
\usepackage{bm}
\usepackage{mathrsfs}
\usepackage{extarrows}

% ---------- Enumeration ----------
\usepackage[shortlabels]{enumitem}
\setlist[enumerate]{itemsep=0mm}

% ---------- Custom commands (same as assignment.cls) ----------
\newcommand{\bvec}[1]{\boldsymbol{#1}}
\newcommand{\ct}{\mathsf{H}}
\newcommand{\diag}{\mathrm{diag}}
\newcommand{\im}{\mathrm{Im}}
\newcommand{\rank}{\mathrm{rank}}
\newcommand{\Span}{\mathrm{span}}

\begin{document}

\title{高等代数 II\quad 习题集}
\author{evilayuria \quad 2025999999}
\date{\today}
\maketitle

\tableofcontents
\newpage

% ============================================
\section{二次型与对称矩阵、正定性}
\section*{Problem 1}

给定二次型
\[
Q \left(x_1, x_2, x_3\right) = 2x_1^2 + 3x_2^2 + x_3^2 + 4x_1x_2 - 2x_1x_3 + 6x_2x_3.
\]
写出其对应的对称矩阵.

\section*{Problem 2}

判断下列二次型的正定性：

\begin{enumerate}[]
    \item $Q_1 \left(x, y\right) = x^2 + 2xy + 2y^2, $
    \item $Q_2 \left(x, y, z\right) = x^2 + y^2 + z^2 - xy - xz - yz.$
\end{enumerate}

\section*{Problem 3}

设 $f(\boldsymbol{x}) = \boldsymbol{x}^\top A \boldsymbol{x}$，定义双线性形
\[
g\left(\boldsymbol{x}, \boldsymbol{y}\right) = \frac{1}{2} \left[f \left(\boldsymbol{x} + \boldsymbol{y}\right) - f\left(\boldsymbol{x}\right) - f\left(\boldsymbol{y}\right)\right].
\]
若
\[
A =
\begin{bmatrix}
2 & 1 \\
1 & 3
\end{bmatrix},
\]
求 $g \left(\boldsymbol{e}_1, \boldsymbol{e}_2\right)$，其中 $\boldsymbol{e}_1 = \left(1, 0\right)^\top, \boldsymbol{e}_2 = \left(0, 1\right)^\top$.

\section*{Problem 4}

将二次型
\[
Q \left(x, y\right) = 2x^2 + 4xy + 5y^2
\]
化为标准形，并写出所用的可逆变换矩阵 $P$.


\newpage
\section{二次型的标准形、惯性指数与正定判别}
\section*{Problem 1}

将
\[
Q\left(\boldsymbol{x}\right) = x_1^2 + 2x_2^2 + 3x_3^2 + 2x_1x_2 + 4x_1x_3 + 2x_2x_3
\]
化为标准形.

\section*{Problem 2}

设
\[
Q\left(x_1, x_2, x_3\right) = ax_1^2 + bx_2^2 + ax_3^2 + 2cx_1x_3,
\]
求 $a, b, c$ 满足什么条件时 $Q$ 为正定二次型.

\section*{Problem 3}

设
\[
Q \left(\boldsymbol{x}\right) = x_1^2 + 2x_2^2 + 3x_3^2 - 2x_1x_2 + 4x_2x_3
\]
求其正负惯性指数.

\section*{Problem 4}

设
\[
A =
\begin{bmatrix}
    2 & -1 & 1 \\
    -1 & 2 & -1 \\
    1 & -1 & 0 \\
\end{bmatrix}
\]
\begin{enumerate}
    \item 构造 $P$ 使 $P^\top A P$ 为标准形；
    \item 求惯性指数；
    \item 判断是否正定.
\end{enumerate}


\newpage
\section{双线性型、Gram 矩阵与对称双线性型}
\section*{Problem 1}

设 $g : \mathbb{R}^2 \times \mathbb{R}^2 \to \mathbb{R}$ 定义为
\[
g \left(\boldsymbol{x}, \boldsymbol{y}\right) = 2 x_1 y_1 + 3 x_1 y_2 - x_2 y_1 + 4 x_2 y_2,
\]
其中 $\boldsymbol{x} = \left(x_1, x_2\right), \boldsymbol{y} = \left(y_1, y_2\right)$.
\begin{enumerate}
    \item 写出 $g$ 在标准基下的 Gram 矩阵 $G$；
    \item 判断 $g$ 是否为对称双线性型.
\end{enumerate}

\section*{Problem 2}

设 $V$ 是 2 维线性空间，基为 $\left\{\boldsymbol{e}_1, \boldsymbol{e}_2\right\}$，定义双线性型 $g$ 满足
\[
\begin{aligned}
g \left(\boldsymbol{e}_1, \boldsymbol{e}_1\right) = 1, \\
g \left(\boldsymbol{e}_1, \boldsymbol{e}_2\right) = 2, \\
g \left(\boldsymbol{e}_2, \boldsymbol{e}_1\right) = 3, \\
g \left(\boldsymbol{e}_2, \boldsymbol{e}_2\right) = 4. \\
\end{aligned}
\]
\begin{enumerate}
    \item 写出 $g$ 在基 $\left\{\boldsymbol{e}_1, \boldsymbol{e}_2\right\}$ 下的 Gram 矩阵；
    \item 若新基 $\left\{\boldsymbol{f}_1, \boldsymbol{f}_2\right\}$ 满足 $\boldsymbol{f}_1 = \boldsymbol{e}_1 + \boldsymbol{e}_2, \boldsymbol{f}_2 = \boldsymbol{e}_1 - \boldsymbol{e}_2$，求 $g$ 在新基下的 Gram 矩阵.
\end{enumerate}

\section*{Problem 3}

设 $g \left(\boldsymbol{x}, \boldsymbol{y}\right) = \boldsymbol{x}^\top A \boldsymbol{y}$ 是 $\mathbb{R}^n$ 上的双线性型，$A$ 为 $n \times n$ 矩阵. 证明：$g$ 是对称双线性型，当且仅当 $A$ 是对称矩阵.


\newpage
\section{线性映射与线性变换、核与像、矩阵表示}
\section*{Problem 1}

定义映射 $\mathcal{T}: \mathbb{R}^3 \to \mathbb{R}^2$ 为
\[
T \left(x_1, x_2, x_3\right) = \left(x_1 + 2x_2, 3x_2 - x_3\right)^\top.
\]
\begin{enumerate}
    \item 证明：$\mathcal{T}$ 是线性映射；
    \item 求 $\ker \left(\mathcal{T}\right)$ 和 $\im \left(\mathcal{T}\right)$ 的维数与一组基.
\end{enumerate}

\section*{Problem 2}

设 $V$ 是 3 维线性空间，基为 $\left\{\bvec{v}_1, \bvec{v}_2, \bvec{v}_3\right\}$，线性变换 $\mathcal{A}: V \to V$ 满足
\[
\begin{aligned}
    \mathcal{A} \left(\bvec{v}_1\right) = \bvec{v}_1 + \bvec{v}_2, \\
    \mathcal{A} \left(\bvec{v}_2\right) = \bvec{v}_2 + \bvec{v}_3, \\
    \mathcal{A} \left(\bvec{v}_3\right) = \bvec{v}_1 + \bvec{v}_3. \\
\end{aligned}
\]
\begin{enumerate}
    \item 写出 $\mathcal{A}$ 在基 $\left\{\bvec{v}_1, \bvec{v}_2, \bvec{v}_3\right\}$ 下的矩阵表示；
    \item 求 $\mathcal{A}$ 的像空间 $\im \left(\mathcal{A}\right)$ 和核空间 $\ker \left(\mathcal{A}\right)$ 的基.
\end{enumerate}

\section*{Problem 3}

设 $\mathcal{A}: \mathbb{R}^2 \to \mathbb{R}^2$ 是线性变换，在标准基下的矩阵为
\[
A = \begin{bmatrix}
    1 & 2 \\
    0 & 3 \\
\end{bmatrix}.
\]
\begin{enumerate}
    \item 求 $\mathcal{A}$ 的像空间 $\im \left(\mathcal{A}\right)$ 和核空间 $\ker \left(\mathcal{A}\right)$；
    \item 判断 $\mathcal{A}$ 是否为单射、满射.
\end{enumerate}


\newpage
\section{秩--零化度定理、基变换与相似矩阵、可对角化判定}
\section*{Problem 1}

给定线性映射 $T : \mathbb{R}^3 \to \mathbb{R}^3$，其在标准基下对应的矩阵为
\[
A =
\begin{bmatrix}
1 & 2 & 3 \\
2 & 4 & 6 \\
1 & 2 & 3
\end{bmatrix}.
\]
请计算矩阵 $A$ 的秩 $\rank \left(A\right)$，并依此求出 $T$ 的像空间维数 $\dim \im \left(\mathcal{T}\right)$ 与核空间维数 $\dim \ker\left(\mathcal{T}\right)$，最后说明其计算结果是否满足秩零化度定理.

\section*{Problem 2}

设线性变换 $T : \mathbb{R}^2 \to \mathbb{R}^2$ 定义为
\[
T\begin{pmatrix} x_1 \\ x_2 \end{pmatrix}
=
\begin{pmatrix}
2x_1 + x_2 \\
x_1 + 2x_2
\end{pmatrix}.
\]
\begin{enumerate}
\item 求 $T$ 在标准基 $\bvec{e}_1 = \begin{pmatrix}1\\0\end{pmatrix},\;
\bvec{e}_2 = \begin{pmatrix}0\\1\end{pmatrix}$ 下对应的矩阵 $A$；
\item 给定新基（已固定顺序）
\[
\tilde{\bvec{e}}_1 = \begin{pmatrix}1\\1\end{pmatrix}, \quad
\tilde{\bvec{e}}_2 = \begin{pmatrix}-1\\1\end{pmatrix},
\]
求从标准基到新基的过渡矩阵 $P$，并计算 $T$ 在该新基下对应的矩阵 $B$.
\end{enumerate}

\section*{Problem 3}

设矩阵
\[
A =
\begin{bmatrix}
4 & -2 \\
1 & 1
\end{bmatrix}.
\]
求矩阵 $A$ 的所有特征值.对于每一个特征值，求其对应的特征子空间的维数（即几何重数），并据此判断矩阵 $A$ 是否可对角化.


\newpage
\section{子空间的直和、维数公式与对称/反对称分解}
\section*{Problem 1}

在 $\mathbb{R}^3$ 中，设子空间 $W_1$ 由方程 $x+y-z=0$ 确定，子空间 $W_2$ 由方程 $x-y=0$ 确定.请分别计算子空间 $W_1$、$W_2$ 以及交空间 $W_1 \cap W_2$ 的维数，并利用维数公式求和空间 $W_1 + W_2$ 的维数.

\section*{Problem 2}

设 $\mathbb{R}^4$ 中的两个子空间分别为
\[
\begin{aligned}
    U &= \Span\left\{
    \left(1, 0, 0, 0\right)^\top,
    \left(0, 1, 0, 0\right)^\top
    \right\}, \\
    V &= \Span\left\{
    \left(1, 1, 1, 1\right)^\top,
    \left(0, 0, 1, 1\right)^\top
    \right\}.
\end{aligned}
\]
请通过计算维数 $\dim \left(U + V\right)$，来判断 $U + V$ 是否为直和（即判断是否满足 $U \oplus V$）.

\section*{Problem 3}

已知 $n$ 阶方阵空间 $V = \mathbb{R}^{n \times n}$ 可以分解为对称矩阵子空间 $W_1$ 与反对称矩阵子空间 $W_2$ 的直和（即 $V = W_1 \oplus W_2$）.给定矩阵
\[
A =
\begin{bmatrix}
1 & 2 & 4 \\
4 & 5 & 6 \\
2 & 8 & 9
\end{bmatrix},
\]
请利用直和分解的性质，求出唯一确定的对称矩阵 $S \in W_1$ 和反对称矩阵 $K \in W_2$，使得 $A = S + K$.


\newpage
\section{不变子空间与分块矩阵}
\section*{Problem 1}

设线性变换 $\mathcal{A} : \mathbb{R}^3 \to \mathbb{R}^3$ 的矩阵为
\[
A =
\begin{bmatrix}
2 & 1 & 0 \\
0 & 2 & 0 \\
0 & 0 & 3
\end{bmatrix}.
\]
判断下列子空间是否为 $\mathcal{A}$ 的不变子空间，并说明理由：
\[
W_1 = \Span\{(1,0,0)^\top\},\quad
W_2 = \Span\{(1,0,0)^\top,(0,1,0)^\top\},\quad
W_3 = \Span\{(1,1,0)^\top\}.
\]

\section*{Problem 2}

设
\[
A =
\begin{bmatrix}
A_1 & A_2 \\
O & A_3
\end{bmatrix}
\]
为 $n$ 阶分块矩阵，其中 $A_1$ 为 $r$ 阶方阵。证明子空间
\[
W_1 = \left\{\left(\alpha_1,\ldots,\alpha_r,0,\ldots,0\right)^\top : \alpha_i \in \mathbb{R}\right\}
\]
是 $A$ 的不变子空间；并证明
\[
W_2 = \left\{\left(0,\ldots,0,\beta_{r+1},\ldots,\beta_n\right)^\top : \beta_j \in \mathbb{R}\right\}
\]
是 $A$ 的不变子空间当且仅当 $A_2 = O$。


\newpage
\section{赋范向量空间、Banach 不动点定理与算子范数}
\section*{Problem 1}

在 $\mathbb{R}^2$ 上定义无穷范数
\[
\|(x,y)\|_\infty = \max\{|x|,|y|\}.
\]
设映射 $\mathcal{T} : \mathbb{R}^2 \to \mathbb{R}^2$ 定义为
\[
\mathcal{T} \left(x,y\right) = \left(\frac14 x + \frac12 y + 1,\; \frac13 x + \frac16 y - 2\right)^\top.
\]
证明 $\mathcal{T}$ 是压缩映射，并求其唯一不动点。

\section*{Problem 2}

设 $\mathcal{T} : \mathbb{R} \to \mathbb{R}$ 定义为
\[
\mathcal{T}(x) = \frac12 \cos x.
\]
证明方程
\[
x = \frac12 \cos x
\]
有唯一实根，并说明迭代
\[
x_{n+1} = \mathcal{T}(x_n)
\]
对任意初值 $x_0 \in \mathbb{R}$ 都收敛到该实根。

\section*{Problem 3}

设
\[
A =
\begin{bmatrix}
1 & 2 \\
2 & 4
\end{bmatrix}.
\]
在 $\mathbb{R}^2$ 中取 Euclid 范数 $\|\cdot\|_2$，求算子范数 $\|A\|_2$。

\section*{Problem 4}

考虑初值问题
\[
y^\prime \left(x\right) = y\left(x\right),\quad y\left(0\right) = 1.
\]
将其改写为积分方程
\[
y\left(x\right) = 1 + \int_0^x y\left(t\right)\, \mathrm{d} t.
\]
在区间 $[0,a]$ 上的连续函数空间 $C\left[0,a\right]$ 中取范数
\[
\left\|y\right\|_\infty = \max_{x \in \left[0,a\right]} \left|y\left(x\right)\right|.
\]
定义算子
\[
\left(Ay\right)\left(x\right) = 1 + \int_0^x y\left(t\right)\, \mathrm{d} t.
\]
证明当 $0 < a < 1$ 时，$A$ 为压缩映射，从而该初值问题在 $C\left[0,a\right]$ 中有唯一解。


\newpage
\section{主成分分析（PCA）：特征值分解与低秩近似}
\section*{Problem 1}

已知矩阵 \[
A = \begin{bmatrix}
1 & 2 & 3 \\
2 & 4 & 6 \\
1 & 2 & 3 \\
\end{bmatrix}.
\]
该矩阵的各行之间存在明显的倍数关系，因此可以用低秩形式表示。

\begin{enumerate}
    \item 将矩阵 $A$ 写成外积形式 \[A = \boldsymbol{u}\boldsymbol{v}^\top.\]
    \item 求矩阵 $A$ 的秩。
    \item 若直接存储矩阵 $A$ 需要存储全部 $3 \times 3$ 个元素，而采用外积形式时只需存储向量 $\boldsymbol{u}$ 和 $\boldsymbol{v}$ 的全部分量，求：
    \begin{enumerate}
        \item 原矩阵的存储量；
        \item 外积表示下的存储量；
        \item 传输效率提升倍数（定义为压缩前存储量 / 压缩后存储量）。
    \end{enumerate}
\end{enumerate}

\section*{Problem 2}

设对称矩阵 \[
A = \begin{bmatrix}
    3 & 1 \\
    1 & 3 \\
\end{bmatrix}.
\]
利用特征值分解研究该矩阵的低秩近似。

\begin{enumerate}
    \item 求矩阵 $A$ 的全部特征值及对应的单位特征向量；
    \item 写出矩阵 $A$ 的特征值分解形式；
    \item 取最大特征值所对应的一项，写出矩阵 $A$ 的最佳秩 $1$ 近似矩阵 $\hat{A}$。
\end{enumerate}

\section*{Problem 3}

设对称矩阵 \[
A = \begin{bmatrix}
    2 & 1 & 0 \\
    1 & 2 & 0 \\
    0 & 0 & 1 \\
\end{bmatrix}
\]
\begin{enumerate}
    \item 求矩阵 $A$ 的全部特征值及对应的一组单位特征向量；
    \item 按特征值从大到小排序，并写出前三个主成分方向；
    \item 设 $\hat{A}_k$ 为保留前 $k$ 个主成分得到的近似矩阵，判断当误差要求 \[\left\|A - \hat{A}_k\right\| \le 1\] 时，至少需要保留几个主成分。
\end{enumerate}


\newpage
\section{PCA 与协方差矩阵、方差最大化与主成分方向}
\section*{Problem 1}

给定数据矩阵 \[
X = \begin{bmatrix}
    1 & 2 & 3 & 4 \\
    2 & 1 & 0 & 3 \\
\end{bmatrix}
\]
\begin{enumerate}
    \item 求两个特征的样本均值，并写出中心化后的数据矩阵 $\tilde{X}$；
    \item 求 $\tilde{X}\tilde{X}^\top$ 的第一主成分方向。
\end{enumerate}

\section*{Problem 2}

设随机向量 $\boldsymbol{x}$ 的协方差矩阵为 \[
\Sigma = \begin{bmatrix}
    5 & 2 \\
    2 & 2 \\
\end{bmatrix}.
\]
\begin{enumerate}
    \item 证明 \[\mathrm{Var} \left(\boldsymbol{a}^\top \boldsymbol{x} \right) = \boldsymbol{a}^\top \Sigma \boldsymbol{a}.\]
    \item 在约束 $\left\|\boldsymbol{a}\right\| = 1$ 下，求使方差最大的方向。
    \item 写出最大方差，并说明其所对应方向的意义。
\end{enumerate}

\section*{Problem 3}

设协方差矩阵 \[\Sigma = \begin{bmatrix}
    4 & 1 & 0 \\
    1 & 4 & 0 \\
    0 & 0 & 1 \\
\end{bmatrix}.\]
\begin{enumerate}
    \item 求 $\Sigma$ 的全部特征值及对应的一组单位特征向量。
    \item 求前 1 个、前 2 个主成分捕获的总方差。
    \item 计算前 1 个、前 2 个主成分的方差解释比例，并说明其意义。
\end{enumerate}


\newpage
\section{奇异值分解（SVD）、奇异值与低秩近似}
\section*{Problem 1}

设矩阵 \[
A = \begin{bmatrix}
    3 & 0 \\
    0 & 2 \\
    0 & 0 \\
\end{bmatrix}.
\]
\begin{enumerate}
    \item 计算 $A^\top A$，求其特征值与单位特征向量，写出 $A$ 的奇异值 $\sigma_1 > \sigma_2$ 及右奇异向量 $\boldsymbol{v}_1, \boldsymbol{v}_2$。
    \item 利用公式 \[\boldsymbol{u}_i = \frac{A\boldsymbol{v}_i}{\sigma_i},\] 求左奇异向量 $\boldsymbol{u}_1, \boldsymbol{u}_2$，并将其补全为 $\mathbb{R}^3$ 的标准正交基。
    \item 写出 $A$ 的奇异值分解 $A = U \Sigma V^\top$。
\end{enumerate}

\section*{Problem 2}

设矩阵
\[A =
\begin{bmatrix}
    3 & 0 \\
    0 & 1 \\
\end{bmatrix}.
\]
\begin{enumerate}
    \item 直接写出 $A$ 的 SVD（该矩阵已是对角形式，奇异值即对角元）。
    \item 利用谱范数定义 \[\|A\|_2 = \sup_{x \neq 0} \frac{\|Ax\|_2}{\|x\|_2},\]证明 $\|A\|_2 = \sigma_1 = 3$。
    \item 写出 $A$ 的秩 - 1 近似 $A_1$，计算 $\|A - A_1\|_2$。
\end{enumerate}

\section*{Problem 3}

设矩阵
\[
A =
\begin{bmatrix}
    3 & 0 & 0 \\
    0 & 2 & 0 \\
    0 & 0 & 1 \\
\end{bmatrix},
\]
其奇异值为 \[\sigma_1 = 3, \quad \sigma_2 = 2, \quad \sigma_3 = 1.\]
\begin{enumerate}
    \item 写出 $A$ 的秩 - 1 近似 $A_1$ 与秩 - 2 近似 $A_2$，分别计算近似误差 $\|A - A_1\|_2$ 与 $\|A - A_2\|_2$。
    \item 由 Eckart – Young 定理，$A_1$ 是 $A$ 的最优秩 - 1 近似。若另取秩 - 1 矩阵 \[B = 3 \boldsymbol{e}_1 \boldsymbol{e}_2^\top,\]其中 $\boldsymbol{e}_1, \boldsymbol{e}_2$ 为标准基向量，计算 $\|A - B\|_2$ 并验证其不小于 $\|A - A_1\|_2$。
\end{enumerate}


\newpage
\section{SVD 与四个基本子空间、极分解与数值稳定性}
\section*{Problem 1}

设矩阵
\[
A = \begin{bmatrix}
    1 & 2 \\
    2 & 4 \\
\end{bmatrix}.
\]
\begin{enumerate}
    \item 求 $A$ 的秩，以及列空间 $C \left(A\right)$ 与零空间 $N \left(A\right)$ 各自的一组基。
    \item 对 $A$ 进行 SVD，在该矩阵的 SVD 中，说明非零奇异值对应的左、右奇异向量分别张成哪个基本子空间；补齐后的其余左、右奇异向量又分别张成哪个基本子空间。
    \item 将向量 $\boldsymbol{x} = \left(3, 1\right)^\top$ 分解为行空间分量与零空间分量之和，即写出 $\boldsymbol{x} = \boldsymbol{x}_{R} + \boldsymbol{x}_{N}$。
\end{enumerate}

\section*{Problem 2}

设矩阵
\[
A =
\begin{bmatrix}
    0 & -1 \\
    2 & 0 \\
\end{bmatrix}.
\]
\begin{enumerate}
    \item 对 $A$ 进行 SVD，求 $U, \Sigma, V$。
    \item 由 SVD 写出极分解 $A = QS$，其中 $Q = UV^\top$，$S = V \Sigma V^\top$，明确给出 $Q$ 和 $S$ 的数值。
    \item 验证 $S$ 是对称半正定矩阵，并说明 $S$ 的特征值与 $A$ 的奇异值有何关系。
\end{enumerate}

\section*{Problem 3}

设矩阵
\[
A = \begin{bmatrix}
    0 & 2 \\
    0 & 0 \\
\end{bmatrix}, \quad A_\varepsilon = \begin{bmatrix}
    0 & 2 \\
    \varepsilon & 0 \\
\end{bmatrix}, \quad \varepsilon > 0.
\]
\begin{enumerate}
    \item 计算 $A$ 的全部特征值；再计算 $A_\varepsilon$ 的特征值，说明当 $\varepsilon$ 很小时特征值的变化情况。
    \item 分别计算 $A^\top A$ 与 $A_\varepsilon^\top A_\varepsilon$ 的特征值，写出 $A$ 与 $A_\varepsilon$ 各自的奇异值，比较二者的差异。
    \item 综合以上结果，说明奇异值相比特征值具有更好数值稳定性的原因。
\end{enumerate}


\newpage
\section{复矩阵与 Hermite 矩阵、酉对角化与 Fourier 矩阵}
\section*{Problem 1}

设 \[A = \begin{bmatrix}
    3 & 2 - \mathrm{i} \\
    2 + \mathrm{i} & 7 \\
\end{bmatrix}.\]

\begin{enumerate}
    \item 证明 $A$ 是 Hermite 矩阵。
    \item 求 $A$ 的特征值与对应的特征向量。
    \item 求酉矩阵 $U$ 使得 $U^\ct AU$ 为对角矩阵，并写出该对角矩阵。
\end{enumerate}

\section*{Problem 2}

设 $A$ 为 $n$ 阶 Hermite 矩阵。证明：
\begin{enumerate}
    \item $I + \mathrm{i}A$ 可逆。
    \item 矩阵 $U = (I - \mathrm{i}A)(I + \mathrm{i}A)^{-1}$ 是酉矩阵。
\end{enumerate}

\section*{Problem 3}

设 $F_4$ 为 4 阶 Fourier 矩阵。

\begin{enumerate}
    \item 写出 $F_4$ 的具体形式（即写出每个元素）；
    \item 证明 $F_4$ 是酉矩阵；
    \item 计算 $F_4\bm{v}$，其中 $\bm{v} = (1, \mathrm{i}, -1, -\mathrm{i})^\top$。
\end{enumerate}


\newpage
\section{Hermite 正定矩阵、伴随算子与复内积空间}
\section*{Problem 1}

设 \[
A = \begin{bmatrix}
    1 & 1 + \mathrm{i} \\
    1 - \mathrm{i} & 3 \\
\end{bmatrix}.
\]

\begin{enumerate}
    \item 证明 $A$ 是 Hermite 正定矩阵。
    \item 求可逆矩阵 $P$ 使得 $P^\ct A P = I$。
    \item 写出 $A$ 对应的 Hermite 二次型 $H \left(\bm{x}\right) = \bm{x}^\ct A \bm{x}$ 在某组新基下的标准形。
\end{enumerate}

\section*{Problem 2}

设 $M$ 是 $m \times n$ 的复矩阵。证明：

\begin{enumerate}
    \item $M^\ct M$ 是 Hermite 半正定矩阵。
    \item $M^\ct M$ 是正定矩阵当且仅当 $M$ 的列向量线性无关。
    \item 任意 Hermite 正定矩阵 $A$ 都可以写成 $A = B^\ct B$，其中 $B$ 是某可逆矩阵。
\end{enumerate}

\section*{Problem 3}

设 $V$ 为有限维复内积空间（内积为 Hermite 内积），$\mathcal{A}, \mathcal{B}$ 为 $V$ 上的线性算子，$\mathcal{A^*}, \mathcal{B^*}$ 分别为其伴随算子。证明：

\begin{enumerate}
    \item $\left(\mathcal{AB}\right)^* = \mathcal{B^*}\mathcal{A^*}.$
    \item 若 $\lambda$ 为 $\mathcal{A}$ 的特征值，则 $\bar\lambda$ 是 $\mathcal{A^*}$ 的特征值。
    \item 若 $\mathcal{A}$ 是 Hermite 算子（即 $\mathcal{A^*} = \mathcal{A}$）且可逆，则 $\mathcal{A}^{-1}$ 也是 Hermite 算子。
\end{enumerate}


\newpage
\section{Fourier 变换与 FFT：蝶形算法、多项式插值与复杂度分析}
\section*{Problem 1}

设 $\omega_n = e^{-2\pi \mathrm{i}/n}$ 为 $n$ 次单位根。请证明：对任意整数
$a, b \in \{0, 1, \ldots, n - 1\}$，
\[
    \sum_{k = 0}^{n - 1} \omega_n^{(a - b)k}
    =
    \begin{cases}
        n, & a = b, \\
        0, & a \neq b.
    \end{cases}
\]
并由此推出 $n$ 阶 Fourier 矩阵
\[
    F_n = \left(\omega_n^{jk}\right)_{j,k = 0}^{n - 1}
\]
满足
\[
    F_n^\ct F_n = n I_n, \qquad F_n^{-1} = \frac{1}{n} F_n^\ct.
\]

\section*{Problem 2}

取 $n = 4$，$\omega_4 = e^{-2\pi \mathrm{i}/4} = -\mathrm{i}$。设输入向量
\[
    \bm{x} = \left(1, 2, 3, 4\right)^\top.
\]

\begin{enumerate}
    \item 直接利用 Fourier 矩阵 $F_4$ 计算 $\bm{y} = F_4 \bm{x}$。
    \item 用 FFT 的蝶形分解（即先分别计算偶数项 $\left(x_0, x_2\right)$ 与奇数项 $\left(x_1, x_3\right)$ 的 2 点 DFT，再用旋转因子合并）重新计算 $\bm{y}$，并核对两种方法结果一致。
\end{enumerate}

\section*{Problem 3}

设次数小于 $n$ 的多项式
\[
    p(z) = x_0 + x_1 z + x_2 z^2 + \cdots + x_{n - 1} z^{n - 1},
\]
记 $\omega_n = e^{-2\pi \mathrm{i}/n}$。

\begin{enumerate}
    \item 证明 $\bm{y} = F_n \bm{x}$ 的第 $j$ 个分量恰为 $y_j = p(\omega_n^j)$，即 DFT 是把多项式从“系数表示”转为“在 $n$ 次单位根处的点值表示”。
    \item 设 $n = 2m$。把 $p(z)$ 按偶、奇次项拆分为
    \[
        p(z) = p_e(z^2) + z p_o(z^2),
    \]
    利用 $\omega_{2m}^2 = \omega_m$ 与 $\omega_{2m}^m = -1$ 推导 FFT 的蝶形合并公式
    \[
        y_j = p_e(\omega_m^j) + \omega_{2m}^j p_o(\omega_m^j), \qquad
        y_{j + m} = p_e(\omega_m^j) - \omega_{2m}^j p_o(\omega_m^j),
    \]
    其中 $j = 0, 1, \ldots, m - 1$。
    \item 写出递推式 $T(n) = 2T(n/2) + O(n)$ 的解，并说明 FFT 相比直接计算 DFT 在复杂度上的优势。
\end{enumerate}


\newpage
\section{矩阵微积分：向量/矩阵求导、迹的导数与行列式求导}
\section*{Problem 1}

设 $\bm{a}, \bm{b} \in \mathbb{R}^n$ 为常向量，$A \in \mathbb{R}^{n \times n}$ 为常矩阵（不必对称），$\bm{x} \in \mathbb{R}^n$。计算下列各式对 $\bm{x}$ 的导数：

\begin{enumerate}
    \item $f_1(\bm{x}) = \bm{a}^\top \bm{x}$；
    \item $f_2(\bm{x}) = \bm{x}^\top A \bm{x}$（$A$ 不必对称）；
    \item $f_3(\bm{x}) = \left(\bm{a}^\top \bm{x}\right)\left(\bm{b}^\top \bm{x}\right)$。
\end{enumerate}

\section*{Problem 2}

设 $X \in \mathbb{R}^{n \times n}$ 为可逆矩阵，$A, B$ 为常矩阵，使下列各式有意义。利用矩阵求导的链式法则及行列式求导公式，计算：

\begin{enumerate}
    \item $\displaystyle \frac{\partial}{\partial X} \ln \det(X)$（不假设 $X$ 对称）；
    \item $\displaystyle \frac{\partial}{\partial X} \ln \det(AXB)$（设 $AXB$ 可逆）；
    \item 验证当 $X$ 为对称可逆矩阵时，
    \[
        \frac{\partial}{\partial X} \ln \det(X)
        =
        2X^{-1} - \diag(X^{-1}).
    \]
\end{enumerate}

\section*{Problem 3}

设 $X \in \mathbb{R}^{n \times m}$，$A \in \mathbb{R}^{m \times m}$，$B \in \mathbb{R}^{n \times n}$ 均为常矩阵。利用 $\displaystyle \frac{\partial \operatorname{tr}(\cdot)}{\partial X}$ 的求导技巧，证明：

\begin{enumerate}
    \item $\displaystyle \frac{\partial}{\partial X} \operatorname{tr}\left(XAX^\top\right) = X(A + A^\top)$；
    \item $\displaystyle \frac{\partial}{\partial X} \operatorname{tr}\left(X^\top BXA\right) = BXA + B^\top XA^\top$；
    \item 作为 (2) 的应用，设 $B$ 对称且 $X^\top BX$ 可逆，证明
    \[
        \frac{\partial}{\partial X} \ln \det\left(X^\top BX\right)
        =
        2BX\left(X^\top BX\right)^{-1}.
    \]
\end{enumerate}


\newpage
\section{正规方程、Sigmoid 导数与 PCA 的 Lagrange 对偶}
\section*{Problem 1}

在多元线性回归中，假设设计矩阵 $X$ 列满秩。为了极小化最小二乘目标函数
\[
    \left\|X\bm{\beta} - \bm{y}\right\|^2,
\]
请推导参数 $\bm{\beta}$ 的最优解所满足的正规方程（Normal Equation）。

\section*{Problem 2}

在单层神经网络
\[
    \bm{z} = \sigma\left(W\bm{x} + \bm{b}\right)
\]
中，设激活函数为 Sigmoid 函数
\[
    \sigma(t) = \frac{1}{1 + e^{-t}}.
\]
请推导导数 $\sigma'(t)$，并将其恒等变形为仅包含 $\sigma(t)$ 的表达式。

\section*{Problem 3}

在主成分分析（PCA）中，求解半正定矩阵 $A^\top A$ 的最大特征值问题等价于求解约束优化问题
\[
    \max_{\|\bm{x}\| = 1}\bm{x}^\top A^\top A\bm{x}.
\]
请利用拉格朗日（Lagrange）乘子法证明：该优化问题的最大值即为矩阵 $A^\top A$ 的最大特征值 $\lambda_1$。


\newpage
\section{Moore--Penrose 广义逆与最小二范数解}
\section*{Problem 1}

已知矩阵
\[
    A = \begin{bmatrix}
        1 & 1 \\
        1 & -1 \\
        1 & 0
    \end{bmatrix}.
\]
由于矩阵 $A$ 列满秩，其 Moore-Penrose 广义逆 $A^+$ 可以通过公式
\[
    A^+ = \left(A^\top A\right)^{-1}A^\top
\]
来计算。请写出详细的计算过程，求出 $A^+$。

\section*{Problem 2}

假设线性方程组 $A\bm{x} = \bm{b}$ 有解。利用广义逆的通解性质可知，该方程组的解可以统一表示为
\[
    \bm{x} = A^+\bm{b} + \left(I - A^+A\right)\bm{z},
\]
其中 $\bm{z}$ 为任意向量，且 $A^+$ 为 $A$ 的 Moore-Penrose 广义逆。请证明：
\[
    \bm{x}^* = A^+\bm{b}
\]
是该方程组所有解中二范数（模长）最小的解。

\section*{Problem 3}

已知 $A^+$ 是矩阵 $A$ 的 Moore-Penrose 广义逆，满足以下四条定义性质：
\[
    \begin{aligned}
        &\text{(1)}\quad AA^+A = A; &
        &\text{(2)}\quad A^+AA^+ = A^+; \\
        &\text{(3)}\quad \left(AA^+\right)^\top = AA^+; &
        &\text{(4)}\quad \left(A^+A\right)^\top = A^+A.
    \end{aligned}
\]
请利用上述四条性质，证明 $\left(A^+\right)^\top$ 也是 $A^\top$ 的 Moore-Penrose 广义逆，从而得出结论
\[
    \left(A^\top\right)^+ = \left(A^+\right)^\top.
\]


\newpage
\section{人工神经元：前向计算、激活函数与梯度消失}
\section*{Problem 1}

一个人工神经元的输入为
\[
\bm{x} = (1, -2, 3)^\top,\quad
\bm{w} = (2, -1, 0.5)^\top,\quad
b = -1.
\]
记净输入为 $z = \bm{w}^\top \bm{x} + b$。
\begin{enumerate}
    \item 计算该神经元的净输入 $z$。
    \item 若激活函数为阶梯函数
    \[
    h(z) = \begin{cases}
        1, & z \ge 0,\\
        0, & z < 0,
    \end{cases}
    \]
    求神经元输出。
    \item 若激活函数分别为 $\operatorname{ReLU}(z) = \max\{0, z\}$ 和 LeakyReLU，其中
    \[
    \operatorname{LeakyReLU}(z) = \begin{cases}
        z,      & z \ge 0,\\
        0.01z,  & z < 0,
    \end{cases}
    \]
    分别求神经元输出，并说明二者在 $z < 0$ 时的主要区别。
\end{enumerate}

\section*{Problem 2}

Logistic 激活函数定义为
\[
\sigma(z) = \frac{1}{1 + \exp(-z)}.
\]
已知其导数为 $\sigma'(z) = \sigma(z)(1 - \sigma(z))$。
\begin{enumerate}
    \item 计算 $\sigma(0)$ 与 $\sigma'(0)$。
    \item 近似计算 $\sigma(4), \sigma'(4), \sigma(-4), \sigma'(-4)$。可使用 $\exp(-4) \approx 0.0183$, $\exp(4) \approx 54.6$。
    \item 根据以上结果，解释为什么 Logistic 激活函数在 $|z|$ 较大时容易导致梯度消失。
\end{enumerate}

\section*{Problem 3}

考虑一个使用 ReLU 激活函数的神经元
\[
z = \bm{w}^\top \bm{x} + b,\quad a = \operatorname{ReLU}(z),
\]
设训练集中所有样本都满足 $\bm{x} \ge 0$（逐分量非负）。
\begin{enumerate}
    \item 写出 ReLU 的函数表达式及其在 $z \neq 0$ 处的导数。
    \item 设 $\bm{w} = (-1, -2)^\top$, $b = 0$，对任意满足 $\bm{x} \ge 0$ 的样本，说明为什么有 $z \le 0$。
    \item 若某个神经元对所有训练样本都有 $z < 0$，说明其权重为什么难以通过梯度下降更新，并说明 LeakyReLU 为什么可以缓解这一问题。
\end{enumerate}


\newpage
\section{全连接网络：结构、参数量、前向计算与 Softmax 回归}
\section*{Problem 1}

一个前馈全连接神经网络有输入层、一个隐藏层和输出层。输入维度为 $d = 4$，隐藏层神经元个数为 $n_1 = 3$，输出层神经元个数为 $n_2 = 2$。忽略偏置项时，前向传播为
\[
\bm{a}^{(0)} = \bm{x},\quad
\bm{z}^{(1)} = W^{(1)} \bm{a}^{(0)},\quad
\bm{a}^{(1)} = \sigma_1(\bm{z}^{(1)}),\quad
\bm{z}^{(2)} = W^{(2)} \bm{a}^{(1)}.
\]
\begin{enumerate}
    \item 写出 $W^{(1)}$ 与 $W^{(2)}$ 的维度。
    \item 分别计算忽略偏置项时和加入偏置项时的参数总数。
    \item 简要说明为什么全连接前馈网络中"相邻两层之间的神经元全部两两连接"会使参数量随层宽快速增大。
\end{enumerate}

\section*{Problem 2}

考虑一个一层隐藏层的前馈神经网络：
\[
\bm{a}^{(0)} = \bm{x},\quad
\bm{z}^{(1)} = W^{(1)} \bm{x},\quad
\bm{a}^{(1)} = \operatorname{ReLU}(\bm{z}^{(1)}),\quad
\hat{\bm{y}} = W^{(2)} \bm{a}^{(1)},
\]
其中
\[
\bm{x} = \begin{pmatrix} 1 \\ -1 \end{pmatrix},\quad
W^{(1)} = \begin{bmatrix}
    1 & 2 \\
    -1 & 1 \\
    2 & 0
\end{bmatrix},\quad
W^{(2)} = \begin{bmatrix} 1 & -2 & 0.5 \end{bmatrix}.
\]
\begin{enumerate}
    \item 计算隐藏层净输入 $\bm{z}^{(1)}$。
    \item 计算隐藏层输出 $\bm{a}^{(1)}$。
    \item 计算网络输出 $\hat{\bm{y}}$；若真实标签为 $y = 2$，采用平方损失 $\ell = (\hat{y} - y)^2$，求该样本的损失值。
\end{enumerate}

\section*{Problem 3}

某三分类模型最后一层输出的 logits 为
\[
\bm{z} = (1, 2, 0)^\top.
\]
经过 softmax 后得到预测概率
\[
\hat{y}_c = \frac{\exp(z_c)}{\sum_{j=1}^{3} \exp(z_j)},\quad c = 1, 2, 3.
\]
设真实标签为第 2 类，即 $\bm{y} = (0, 1, 0)^\top$。
\begin{enumerate}
    \item 写出 $\hat{\bm{y}} = \operatorname{softmax}(\bm{z})$ 的各个分量表达式。
    \item 取 $\exp(1) \approx 2.718$, $\exp(2) \approx 7.389$, $\exp(0) = 1$，近似计算 $\hat{\bm{y}}$。
    \item 计算交叉熵损失
    \[
    \ell(\bm{y}, \hat{\bm{y}}) = -\bm{y}^\top \log \hat{\bm{y}}.
    \]
    并记 $\frac{\partial \ell}{\partial \bm{z}} = \hat{\bm{y}} - \bm{y}$，写出该样本对 logits 的梯度。
\end{enumerate}


\newpage
\section{ReLU 激活模式、仿射映射、卷积参数量与 Softmax 梯度}
\section*{Problem 1}

在全连接神经网络中，ReLU 激活函数会将输入空间划分为若干个区域。在每一个固定的激活区域内，网络实际上退化为一个仿射映射。设一个两层全连接神经网络定义为
\[
\bm{z}^{(1)} = W_1 \bm{x} + \bm{b}_1,\quad
\bm{h} = \operatorname{ReLU}(\bm{z}^{(1)}),\quad
\bm{z}^{(2)} = W_2 \bm{h} + \bm{b}_2,
\]
其中
\[
W_1 = \begin{bmatrix}
    1 & -1 \\
    2 & 1 \\
    -1 & 2
\end{bmatrix},\quad
\bm{b}_1 = \begin{pmatrix} 0 \\ -1 \\ 1 \end{pmatrix},\quad
W_2 = \begin{bmatrix} 1 & 0 & 2 \\ -1 & 1 & 0 \end{bmatrix},\quad
\bm{b}_2 = \begin{pmatrix} 1 \\ 0 \end{pmatrix},
\]
对输入
\[
\bm{x} = \begin{pmatrix} 2 \\ 1 \end{pmatrix},
\]
完成以下任务：
\begin{enumerate}
    \item 计算 $\bm{z}^{(1)}$、$\bm{h}$、$\bm{z}^{(2)}$。
    \item 判断该输入对应的 ReLU 激活模式。
    \item 在该激活模式保持不变的区域内，写出网络对应的仿射映射。
\end{enumerate}

\section*{Problem 2}

全连接层的一个主要问题是参数量随输入维度迅速增长。卷积层通过局部连接和权重共享显著减少参数量。设输入是一张大小为 $32 \times 32$ 的 RGB 图像，即输入通道数为 $3$。

现在比较以下两种结构：
\begin{itemize}
    \item 结构 A：将图像展平成向量后，接一个含有 $1024$ 个神经元的全连接层。
    \item 结构 B：使用 $16$ 个 $3 \times 3$ 的卷积核，步长为 $1$，无零填充。
\end{itemize}
假设所有层都含有偏置项。完成以下任务：
\begin{enumerate}
    \item 计算结构 A 的参数量。
    \item 计算结构 B 的参数量。
    \item 计算结构 B 的输出特征图大小。
    \item 若将结构 B 的输出展平后接一个 $10$ 类分类器，计算该分类器的参数量。
\end{enumerate}

\section*{Problem 3}

在多分类任务中，神经网络通常先输出 logits，再通过 softmax 转化为类别概率。设一个三分类模型的 logits 由
\[
\bm{z} = W \bm{x} + \bm{b}
\]
给出，其中
\[
W = \begin{bmatrix}
    1 & 0 \\
    0 & 1 \\
    -1 & 0
\end{bmatrix},\quad
\bm{b} = \begin{pmatrix} 0 \\ 0 \\ 0 \end{pmatrix},\quad
\bm{x} = \begin{pmatrix} 1 \\ 0 \end{pmatrix},
\]
真实标签为第一类，即
\[
\bm{y} = \begin{pmatrix} 1 \\ 0 \\ 0 \end{pmatrix}.
\]
softmax 输出定义为
\[
p_i = \frac{\exp(z_i)}{\sum_{j=1}^{3} \exp(z_j)},
\]
交叉熵损失函数为
\[
L = -\sum_{i=1}^{3} y_i \log p_i.
\]
完成以下任务：
\begin{enumerate}
    \item 计算 $\bm{z}$ 和 $\bm{p}$。
    \item 计算损失 $L$。
    \item 利用 $\nabla_{\bm{z}} L = \bm{p} - \bm{y}$，计算 $\nabla_W L$ 和 $\nabla_{\bm{b}} L$。
\end{enumerate}


\newpage
\section{反向传播、自动微分与一维卷积的矩阵形式}
\section*{Problem 1}

反向传播算法的核心思想是先计算高层误差，再将误差逐层传回低层。设一个两层线性网络为
\[
\bm{h} = W_1 \bm{x},\quad \hat{\bm{y}} = W_2 \bm{h},
\]
其中
\[
W_1 = \begin{bmatrix}
    1 & 2 \\
    0 & -1
\end{bmatrix},\quad
W_2 = \begin{bmatrix}
    1 & 0 \\
    2 & -1
\end{bmatrix},\quad
\bm{x} = \begin{pmatrix} 1 \\ -1 \end{pmatrix},\quad
\bm{y} = \begin{pmatrix} 0 \\ 1 \end{pmatrix}.
\]
损失函数为
\[
L = \frac{1}{2} \|\hat{\bm{y}} - \bm{y}\|_2^2.
\]
完成以下任务：
\begin{enumerate}
    \item 计算 $\bm{h}$、$\hat{\bm{y}}$、$\bm{e} = \hat{\bm{y}} - \bm{y}$。
    \item 计算 $\nabla_{W_2} L$。
    \item 计算 $\nabla_{W_1} L$。
\end{enumerate}

\section*{Problem 2}

自动微分可以看作是系统地应用链式法则。设计算图为
\[
L = \bm{c}^\top \bm{v},\quad
\bm{v} = \tanh(\bm{u}),\quad
\bm{u} = A \bm{x} + \bm{b},
\]
其中
\[
A = \begin{bmatrix}
    1 & -1 \\
    0 & 2
\end{bmatrix},\quad
\bm{x} = \begin{pmatrix} 1 \\ 0 \end{pmatrix},\quad
\bm{b} = \begin{pmatrix} -1 \\ 0 \end{pmatrix},\quad
\bm{c} = \begin{pmatrix} 2 \\ -1 \end{pmatrix}.
\]
这里 $\tanh$ 逐坐标作用于向量。完成以下任务：
\begin{enumerate}
    \item 进行前向计算，求 $\bm{u}$、$\bm{v}$、$L$。
    \item 使用反向模式自动微分思想，求 $\nabla_{\bm{x}} L$、$\nabla_{A} L$、$\nabla_{\bm{b}} L$。
\end{enumerate}

\section*{Problem 3}

一维卷积或互相关可以看作是一个具有特殊稀疏结构的矩阵乘法。设一维输入信号为
\[
\bm{x} = (x_1, x_2, x_3, x_4, x_5)^\top,
\]
滤波器为
\[
\bm{a} = (a_1, a_2, a_3)^\top.
\]
采用步长为 $1$、无零填充的窄互相关运算：
\[
y_i = a_1 x_i + a_2 x_{i+1} + a_3 x_{i+2},\quad i = 1, 2, 3.
\]
完成以下任务：
\begin{enumerate}
    \item 将该运算写成 $\bm{y} = K \bm{x}$ 的矩阵形式。
    \item 令 $\bm{a} = (1, 0, -1)^\top$，$\bm{x} = (2, 1, 3, 0, -2)^\top$，计算 $\bm{y}$。
    \item 若目标输出为 $\bm{t} = (0, 1, 3)^\top$，损失函数为 $L = \frac{1}{2} \|\bm{y} - \bm{t}\|_2^2$，求 $\nabla_{\bm{x}} L$ 和 $\nabla_{\bm{a}} L$。
\end{enumerate}


\end{document}
